\documentclass[11pt, a4paper]{article}

\usepackage[T1]{fontenc}
\usepackage[utf8]{inputenc}
\usepackage{tgpagella}
\usepackage[spanish, es-noshorthands]{babel}
\usepackage{amsmath, amssymb, bm}
\usepackage{booktabs}
\usepackage{tabularx}
\usepackage[table, xcdraw]{xcolor}
\usepackage[most]{tcolorbox}
\usepackage[american]{circuitikz}
\usepackage[margin=2.5cm]{geometry}
\usepackage{fancyhdr}

\pagestyle{fancy}
\fancyhf{}
\setlength{\headheight}{16pt}
\lhead{\textbf{UCB Sede Tarija} | Ing. Mecatrónica}
\chead{}
\rhead{IMT-121 Circuitos Electrónicos I | Ejercicio Guiado S07-1}
\lfoot{Prof: M.Sc. Ing. Bernardo Quiroga Turdera}
\rfoot{Página \thepage}
\renewcommand{\headrulewidth}{0.4pt}

\definecolor{ucbblue}{RGB}{0, 51, 102}

\begin{document}

\begin{tcolorbox}[colback=ucbblue!5!white, colframe=ucbblue, title=\textbf{Hoja de Cálculo Guiada: Regímenes de Carga y Eficiencia}]
    \centering \textbf{Análisis de Barrido Discreto sobre un Equivalente de Thévenin}
\end{tcolorbox}

\section*{Parámetros del Sistema}
Asuma que se ha simplificado una red lineal compleja hasta obtener el siguiente equivalente de Thévenin:
\begin{itemize}
    \item $V_{th} = 12\,\text{V}$
    \item $R_{th} = 3\,\text{k}\Omega$ ($3000\,\Omega$)
\end{itemize}

\section*{Tabla de Barrido Discreto de Carga ($R_L$)}
Complete la siguiente tabla calculando la corriente de carga ($I_L$), voltaje de carga ($V_L$), potencia en la carga ($P_L$), potencia disipada en el modelo interno ($P_{Rth}$) y eficiencia ($\eta$) para cada valor de $R_L$. Utilice unidades estándar o consistentes ($\text{mA}$, $\text{V}$, $\text{mW}$).

\vspace{0.4cm}
\noindent\renewcommand{\arraystretch}{1.8}
\begin{tabularx}{\textwidth}{|*{6}{>{\centering\arraybackslash}X|}}
\hline
\rowcolor{ucbblue!10} 
\textbf{$R_L$ (k$\Omega$)} & \textbf{$I_L$ (mA)} & \textbf{$V_L$ (V)} & \textbf{$P_L$ (mW)} & \textbf{$P_{Rth}$ (mW)} & \textbf{$\eta$ (\%)} \\
\hline
$0.75$ & & & & & \\
\hline
$1.5$  & & & & & \\
\hline
$3.0$  & & & & & \\
\hline
$6.0$  & & & & & \\
\hline
$12.0$ & & & & & \\
\hline
\end{tabularx}

\vspace{0.6cm}
\section*{Interpretación de Resultados}
Responda brevemente a partir de la tabla obtenida:
\begin{enumerate}
    \item ¿En qué caso se encuentra la mayor corriente circulando por el circuito ($I_L$)? \\[0.1cm] \hrulefill
    \item ¿En qué caso se encuentra la mayor caída de voltaje sobre la carga ($V_L$)? \\[0.1cm] \hrulefill
    \item ¿En qué caso se encuentra la mayor potencia útil ($P_L$)? \\[0.1cm] \hrulefill
    \item ¿En qué caso se alcanza la mayor eficiencia energética ($\eta$)? \\[0.1cm] \hrulefill
    \item ¿Por qué la máxima potencia y la máxima eficiencia no ocurren simultáneamente? Explique el compromiso de diseño.
    \vspace{2cm}
    \item Realice un boceto rápido cualitativo (a mano alzada) de $P_L$ en función de $R_L$, marcando explícitamente el punto donde $R_L = R_{th}$.
\end{enumerate}

\end{document}
